\documentclass[12pt]{ximera}
\title{Homework 1: Preference Schedules, Plurality, and Borda Count}
\begin{document}
\begin{abstract}
Sections 1.1--1.3: converting a preference schedule between formats, and tabulating an
election by the Plurality method and by the Borda Count method --- which, as the second
problem shows, can choose different winners.
\end{abstract}
\maketitle

\section*{Sections 1.1--1.3: Ballots, Plurality, and Borda Count}

\begin{problem}
To determine the winner of this year's Hogwarts House Cup, Headmaster Umbridge has
decided to use a preference ballot with her $20$ kittens as voters. Each ballot listed
the candidates Gryffindor, Hufflepuff, Ravenclaw, and Slytherin, and preference was
taken by the order each kitten stepped on the candidate names. These
``printed-names format'' ballots were then organized into the preference schedule
shown in Table 1.

\begin{center}
\begin{tabular}{c|ccccc}
 & 6 & 5 & 5 & 3 & 1 \\\hline
Gryffindor & 3rd & 2nd & 3rd & 2nd & 3rd \\
Hufflepuff & 2nd & 1st & 4th & 1st & 1st \\
Ravenclaw  & 1st & 4th & 1st & 3rd & 4th \\
Slytherin  & 4th & 3rd & 2nd & 4th & 2nd
\end{tabular}

\smallskip
Table 1. Preference schedule, printed-names format.
\end{center}

\begin{prompt}
\textbf{(a)} Rewrite the preference schedule in ``standard format'' --- that is, with
1st, 2nd, 3rd, and 4th listed at the left and candidates written in for each ballot
column. Use G, H, R, and S for the candidates.

\begin{center}
\begin{tabular}{c|ccccc}
Preference & 6 & 5 & 5 & 3 & 1 \\\hline
1st & $\answer{R}$ & $\answer{H}$ & $\answer{R}$ & $\answer{H}$ & $\answer{H}$ \\
2nd & $\answer{H}$ & $\answer{G}$ & $\answer{S}$ & $\answer{G}$ & $\answer{S}$ \\
3rd & $\answer{G}$ & $\answer{S}$ & $\answer{G}$ & $\answer{R}$ & $\answer{G}$ \\
4th & $\answer{S}$ & $\answer{R}$ & $\answer{H}$ & $\answer{S}$ & $\answer{R}$
\end{tabular}
\end{center}

\begin{hint}
Work one column at a time. In the first column, which house was marked 1st? Which was
marked 2nd? Read down the original table looking for the rank you want.
\end{hint}

\begin{feedback}
Each column of the original table lists, for each house, the rank that group of voters
gave it. To convert, read each column and place each house in the row matching its
rank.

For the first column ($6$ voters): Ravenclaw was 1st, Hufflepuff 2nd, Gryffindor 3rd,
Slytherin 4th --- so that column reads R, H, G, S. Repeating for the rest:
\begin{center}
\begin{tabular}{c|ccccc}
Preference & 6 & 5 & 5 & 3 & 1 \\\hline
1st & R & H & R & H & H \\
2nd & H & G & S & G & S \\
3rd & G & S & G & R & G \\
4th & S & R & H & S & R
\end{tabular}
\end{center}
\end{feedback}
\end{prompt}

\begin{prompt}
\textbf{(b)} Using the Plurality method, who wins this year's Hogwarts House Cup?

\begin{multipleChoice}
\choice{Gryffindor}
\choice{Hufflepuff}
\choice[correct]{Ravenclaw}
\choice{Slytherin}
\end{multipleChoice}

How many first-place votes did the winner receive? $\answer{11}$

\begin{feedback}
Plurality counts only first-place votes:
\[
\text{R}: 6+5 = 11, \qquad \text{H}: 5+3+1 = 9, \qquad \text{G}: 0, \qquad \text{S}: 0.
\]
Ravenclaw wins with $11$ of the $20$ first-place votes.

Note that Gryffindor and Slytherin receive \emph{no} first-place votes at all, even
though Gryffindor is ranked 2nd or 3rd on every single ballot. Plurality is blind to
everything below the top line.
\end{feedback}
\end{prompt}
\end{problem}

\begin{problem}
Table 2 gives the preference schedule for an election for Best Pickled Herring of
Trollfjord and surrounding glaciers.

\begin{center}
\begin{tabular}{c|ccccc}
Preference & 27 & 16 & 15 & 13 & 6 \\\hline
1st & A & D & B & B & D \\
2nd & B & E & D & D & C \\
3rd & C & C & C & E & E \\
4th & D & A & E & A & A \\
5th & E & B & A & C & B
\end{tabular}

\smallskip
Table 2. Preference schedule for Best Pickled Herring.
\end{center}

\begin{prompt}
\textbf{(a)} What was the total number of voters?

$\answer{77}$

\begin{feedback}
$27+16+15+13+6=77$ voters.
\end{feedback}
\end{prompt}

\begin{prompt}
\textbf{(b)} Using the Plurality method, which pickled herring wins?

\begin{multipleChoice}
\choice{A}
\choice[correct]{B}
\choice{C}
\choice{D}
\choice{E}
\end{multipleChoice}

\begin{feedback}
First-place votes: A $27$, B $15+13=28$, D $16+6=22$, and C and E none. B wins with
$28$.
\end{feedback}
\end{prompt}

\begin{prompt}
\textbf{(c)} Give the complete ranking by Plurality, with first-place vote counts.

1st: $\answer{B}$ with $\answer{28}$ \quad 2nd: $\answer{A}$ with $\answer{27}$ \quad
3rd: $\answer{D}$ with $\answer{22}$

\smallskip
C and E tie for last, each with $\answer{0}$ first-place votes.

\begin{feedback}
\[ \text{B }28 > \text{A }27 > \text{D }22 > \text{C }0 = \text{E }0. \]
Plurality cannot separate C and E, since neither is anyone's first choice.
\end{feedback}
\end{prompt}

\begin{prompt}
\textbf{(d)} Using the Borda Count method ($5$ points for 1st down to $1$ for 5th),
which pickled herring wins?

\begin{multipleChoice}
\choice{A}
\choice{B}
\choice{C}
\choice[correct]{D}
\choice{E}
\end{multipleChoice}

The winner's Borda score is $\answer{276}$.

\begin{hint}
For each candidate, multiply each column's voter count by the points for that
candidate's position, then add across the columns.
\end{hint}

\begin{feedback}
For D: $27\cdot 2+16\cdot 5+15\cdot 4+13\cdot 4+6\cdot 5=54+80+60+52+30=276$.
The full tally is
\[
\text{D } 276, \quad \text{B } 270, \quad \text{A } 220, \quad \text{C } 211, \quad \text{E } 178,
\]
which totals $1155=77\cdot 15$. \checkmark
\end{feedback}
\end{prompt}

\begin{prompt}
\textbf{(e)} Give the complete ranking by Borda count, with scores.

1st: $\answer{D}$ with $\answer{276}$ \quad 2nd: $\answer{B}$ with $\answer{270}$ \quad
3rd: $\answer{A}$ with $\answer{220}$

\smallskip
4th: $\answer{C}$ with $\answer{211}$ \quad 5th: $\answer{E}$ with $\answer{178}$

\begin{feedback}
\[ \text{D } 276 > \text{B } 270 > \text{A } 220 > \text{C } 211 > \text{E } 178. \]
Here the two methods pick \emph{different winners}. B wins Plurality on first-place
votes, but D is ranked 1st or 2nd on $50$ of the $77$ ballots, and Borda rewards that
broad support enough to overtake B.
\end{feedback}
\end{prompt}
\end{problem}

\end{document}
